\subsection{可微函数的复合}

1.3.1. 假设 F 可赋范。设 \(x_0 \in E\)、\(y_0 \in F\)，U 是 \(x_0\) 的一个邻域，V 是 \(y_0\) 的一个邻域；最后，设 \(f\) 是从 U 到 V 的映射，在 \(x_0\) 处可微，且 \(f(x_0) = y_0\)。若 \(g\) 是从 V 到分离多范数向量空间 G 的映射，在 \(y_0\) 处可微，则从 U 到 G 的映射 \(g \circ f\) 在 \(x_0\) 处可微，且有：

\[
\mathbf {D} (g \circ f) (x _ {0}) = \mathbf {D} g (y _ {0}) \circ \mathbf {D} f (x _ {0}).\tag{1}
\]

若 \(f\) 和 \(g\) 严格可微，则 \(g \circ f\) 也有同样结论。

1.3.2. 设 \(f\) 是在 E 中一点 \(x_0\)\(E\) 的某邻域上定义且取值于 \(F\) 的映射，并在 \(x_0\) 处可微；若 \(u\) 是从 \(F\) 到分离多范数空间 \(G\) 的连续线性映射，则函数 \(u \circ f\) 在 \(x_0\) 处可微，且有：

\[
\mathrm{D} (u \circ f) (x _ {0}) = u \circ \mathrm{D} f (x _ {0}).\tag{2}
\]

1.3.3. 假设 F 是分离多范数向量空间族 \((F_{i})_{i\in I}\) 的乘积；对 I 中每个 i，设 \(f_{i}\) 是在 E 中一点 \(x_{0}\) 的邻域 U 上定义且取值于 F 的映射，并令 \(f = (f_{i})_{i\in I}\)。f 在 \(x_{0}\) 处可微（分别地，严格可微）的充要条件是所有 \(f_{i}\) 都具有该性质；有：

\[
\mathrm{D} _ {h} f (x _ {0}) = (\mathrm{D} _ {h} f _ {i} (x _ {0})) _ {i \in \mathbf {I}} \quad \text { for   every } h \text { in } E.\tag{3}
\]

1.3.4. 若 E = K，则在公式 (1) 和 (3) 中可用 \(Df(x_{0})\) 替代 \(f'(x_{0})\)，用 \(D_{h}f_{i}(x_{0})\) 替代 \(f_{i}'(x_{0})\)。

